\chapter{Introduction}
\label{ch:introduction}

\section{Background}
\label{sec:background}

Road traffic accidents remain one of the leading causes of preventable
death and injury in the modern world. According to the World Health
Organization's \emph{Global Status Report on Road Safety}, approximately
1.19 million people lose their lives in road crashes every year, and
between twenty and fifty million more suffer non-fatal injuries that
often result in long-term disability. Driver-related factors account
for the overwhelming majority of these incidents, with human error
contributing to roughly ninety percent of all recorded crashes.
Historically, attention has been concentrated on the most visible
contributors such as speeding, alcohol intoxication, and fatigue.
However, a growing body of epidemiological and behavioural research
now implicates \emph{acute psychological stress} as a silent and
under-recognised risk factor behind a substantial share of these
events \citep{sharma2012driver}.

Driving is a cognitively demanding task that continuously requires
sustained visual attention, rapid decision making, and fine motor
control. When a driver experiences elevated psychological stress,
several physiological and cognitive changes take place simultaneously:
heart rate variability decreases, breathing becomes shallower, muscle
tension in the shoulders and neck increases, and the sympathetic branch
of the autonomic nervous system begins to dominate cardiovascular
regulation. These changes are accompanied by narrowing of the visual
field (so-called ``tunnel vision''), slower reaction times, reduced
inhibitory control, and a tendency toward impulsive or aggressive
manoeuvres such as tailgating, abrupt lane changes, and failure to
yield. Epidemiological studies have consistently reported that high
levels of self-reported driving stress are associated with a two- to
threefold increase in crash risk, particularly during rush-hour
commuting and long-haul driving.

The sources of driver stress are diverse. They include external
environmental factors such as traffic congestion, adverse weather,
and complex road geometry, as well as internal and personal factors
such as time pressure, interpersonal conflict, workplace deadlines,
and pre-existing anxiety disorders. Commercial drivers -- taxi drivers,
truck drivers, and ride-share operators -- are particularly vulnerable
because their occupational demands combine long hours behind the wheel
with economic pressure and irregular schedules. Younger drivers and
novice drivers constitute another at-risk group because they tend to
experience higher baseline anxiety while simultaneously lacking the
automatised skills that make experienced driving less cognitively
taxing.

Despite the clear public-health significance of this problem,
contemporary vehicles offer remarkably few countermeasures that
target driver stress directly. Modern driver-assistance systems
primarily address perceptual lapses (lane keeping, forward collision
warning, adaptive cruise control) and gross behavioural impairment
(drowsiness detection, alcohol interlocks), but they are largely blind
to the emotional and autonomic state of the person behind the wheel.
This gap motivates the work presented in this report: we aim to close
it using wearable sensors that many drivers already carry on their
wrists, and using machine-learning techniques that can be executed
entirely on-device within the power and memory budget of a smartwatch.

\section{Motivation}
\label{sec:motivation}

Three complementary trends make real-time, software-based driver
stress detection both feasible and timely today.

First, \textbf{consumer wearables have reached sensor-grade maturity}.
Devices such as the Samsung Galaxy Watch, Google Pixel Watch, Fitbit
Sense, and Apple Watch now routinely expose photoplethysmography (PPG),
skin temperature, three-axis acceleration, and -- on selected models --
electrodermal activity (EDA). The same sensor complement that was, only
a decade ago, confined to laboratory-grade equipment such as the Empatica
E4 or the RespiBAN professional chest band is now available in consumer
devices worn for more than twelve hours per day by hundreds of millions
of users worldwide. The WESAD dataset \citep{schmidt2018wesad}, which
underpins this project, was specifically collected using such a wrist
wearable configuration, and therefore its findings translate more or
less directly to mass-market devices.

Second, \textbf{on-device machine learning is no longer a research
curiosity}. Frameworks such as TensorFlow Lite \citep{tflite2023},
ONNX Runtime Mobile, and Core ML provide mature toolchains for running
quantised neural networks on battery-powered hardware. Modern smartwatch
system-on-chips expose hardware-accelerated inference paths that can
execute small recurrent and convolutional models within milliseconds
while consuming only a few milliwatts. Consequently, stress detection
can take place \emph{entirely on the wrist}, with no dependence on
cellular connectivity, no continuous upload of sensitive biometric
data, and no round-trip latency to a cloud endpoint. This is a crucial
property for automotive use cases where network availability is
intermittent and privacy expectations are high.

Third, \textbf{the actuation side of the problem is finally tractable}
thanks to the widespread availability of low-cost Bluetooth Low Energy
(BLE) microcontrollers such as the Arduino Nano 33 BLE. These devices
can drive relays, LEDs, small fans, and infotainment-system inputs,
while speaking a standard GATT profile to any nearby wearable. This
makes it possible to build an \emph{end-to-end} prototype -- from
wearable sensing to cabin actuation -- without bespoke hardware and
without invasive modifications to the vehicle.

A further motivating factor is the \emph{graduated nature} of the
intervention being studied. Unlike alcohol interlocks, which make a
binary decision that either permits or forbids ignition, a stress-aware
cabin control system can respond proportionally. A mildly elevated
stress level can be addressed by dimming interior ambient lighting
and lowering the in-cabin music volume; a moderate level can additionally
switch the LED ambience to a calming blue and introduce soft background
audio; only a critical level triggers an overt visual and auditory
alert. This calibrated, non-punitive form of intervention is more
likely to be accepted by drivers than a coarse, safety-critical
override, and it complements rather than competes with existing
driver-assistance technology.

\section{Project Scope}
\label{sec:scope}

This project is deliberately \textbf{software-focused}. No new hardware
has been designed. All of the sensing hardware used -- wrist PPG, EDA,
skin temperature, and accelerometer -- exists in off-the-shelf consumer
wearables; all of the actuation hardware -- LEDs, PWM-driven cooling
fan, and piezo speaker -- is directly supported by an unmodified Arduino
Nano 33 BLE development board. The contribution of the project lies
entirely in (i) the machine-learning model and its training pipeline,
(ii) the Wear OS Kotlin application that runs the model in real time,
(iii) the BLE protocol layer that conveys stress estimates to the
vehicle, and (iv) the Arduino firmware that translates those estimates
into cabin actions.

The scope of the scientific evaluation is equally bounded. All training
and testing is performed on the publicly available WESAD dataset. No
naturalistic driving data has been collected; we do not claim that the
system has been validated \emph{in vivo} on real drivers, and this
limitation is revisited explicitly in Chapter~\ref{ch:conclusion}.
The system is, however, architected in a way that does not depend on
laboratory assumptions -- it expects streaming sensor data, tolerates
missing samples, and emits its estimates on a fixed thirty-second
cadence suitable for deployment outside the laboratory.

The scope is further bounded on the safety side. The adaptive cabin
control system is explicitly designed to be \emph{advisory}: it alters
lighting, fan speed, and audio profile, but it does not interfere with
steering, braking, or any other primary vehicle control. This keeps
the prototype squarely within the regulatory envelope of an aftermarket
comfort accessory rather than an automated driving component, and it
aligns the project with the graduated-intervention philosophy described
earlier.

\section{Problem Statement}
\label{sec:problem}

Despite extensive prior work on both stress physiology and wearable
biosensing, four concrete gaps remain unaddressed in the literature
and the market:

\begin{enumerate}
  \item \textbf{No real-time multimodal wearable stress classifier
  targeted at drivers.} The most successful stress classifiers
  published on WESAD use offline, PC-class models that are neither
  quantised nor latency-budgeted for a smartwatch, and many of them
  rely on chest-band signals (respiration, ECG) that wrist wearables
  cannot acquire.

  \item \textbf{No adaptive cabin response system driven by
  physiological state.} Commercial vehicles offer static ``comfort
  profiles'' selected manually by the driver; none of them respond
  automatically to the driver's real-time autonomic state, and none
  of them gradate the response across distinct stress levels.

  \item \textbf{Existing solutions are not wearable-constrained.}
  State-of-the-art stress detection pipelines frequently assume access
  to full-resolution raw signals at 64--700~Hz and to tens of megabytes
  of RAM. Consumer smartwatches offer neither. Bridging that gap
  requires aggressive signal downsampling, a small model footprint,
  and quantisation-aware conversion tooling.

  \item \textbf{No integrated end-to-end open system.} To our
  knowledge, no published work combines wearable sensing, on-device
  inference, a documented BLE protocol, and a vehicle-side actuator
  controller in a single openly described prototype.
\end{enumerate}

These gaps together define the problem addressed in this project.

\section{Objectives}
\label{sec:objectives}

To close the gaps identified above, this project pursues the
following five concrete objectives:

\begin{enumerate}
  \item \textbf{Design and train a four-class driver stress classifier}
  (Calm / Mild / Moderate / Critical) based on a bidirectional LSTM
  \citep{graves2005bilstm} with an additive attention layer
  \citep{vaswani2017attention}, using only wrist-accessible signals
  from the WESAD dataset \citep{schmidt2018wesad}.

  \item \textbf{Achieve accuracy above 90\% and macro-F1 above 0.90}
  on a held-out test split while keeping the model size compatible
  with on-device inference (target $\le 80$~KB after TensorFlow Lite
  conversion).

  \item \textbf{Implement a Wear OS Kotlin application} that performs
  real-time sensor collection, windowing, normalisation, and on-device
  TensorFlow Lite inference \citep{tflite2023} with an end-to-end
  latency below fifty milliseconds.

  \item \textbf{Specify and implement a compact BLE protocol}
  (a twelve-byte GATT characteristic) carrying timestamp, stress level,
  and confidence, and connect it to an Arduino Nano 33 BLE that
  implements the vehicle-side finite state machine.

  \item \textbf{Demonstrate a graduated adaptive cabin response}
  consisting of four differentiated cabin profiles (Calm, Mild,
  Moderate, Critical), a waiting state for connection loss, and a
  sixty-second BLE timeout enforced at the Arduino.
\end{enumerate}

Each subsequent chapter of this report addresses one or more of
these objectives directly, and Chapter~\ref{ch:conclusion} revisits
them in order to demonstrate that each has been met by the
implemented system.

\section{Connection to Patent ACN1408}
\label{sec:acn1408}

This work extends a prior invention by the same author, filed as an
Indian Provisional Patent Application under the number ACN1408
\citep{acn1408}. The earlier invention -- commercially referred to
as \emph{AlcoWatch} -- addresses the problem of \textbf{alcohol-related
driving accidents} by using a wrist wearable to estimate blood alcohol
concentration (BAC) from physiological correlates (PPG, EDA, skin
temperature) and by using a relay-controlled Arduino module to block
vehicle ignition when BAC exceeds the legal limit of 0.08~g/dL. In
that system the actuation is binary (ignition ALLOW / ignition BLOCK),
the failure mode is conservative (ignition is blocked by default),
and the model is regressive (it outputs a scalar BAC estimate).

The present project generalises ACN1408 along two orthogonal axes.
First, it moves from a \emph{binary safety interlock} (ignition block)
to a \emph{graduated comfort intervention} (cabin profile). Second,
it moves from \emph{BAC regression} to \emph{stress classification}
with four ordinal classes. The two systems are architecturally
compatible: they reuse the same class of wearable hardware, the same
class of Arduino microcontroller, and a very similar BLE service
structure. Chapter~\ref{ch:architecture} provides a direct side-by-side
comparison, and Chapter~\ref{ch:conclusion} sketches how the two could
be unified into a single multi-purpose wearable safety platform
in future work.
